\documentclass[10pt]{beamer}

\usepackage[utf8]{inputenc}
\usepackage[T1]{fontenc}
\usepackage{pgf}
\usepackage{beamerthemesplit}
\usepackage{graphicx}
\usepackage{algorithm}
\usepackage{algpseudocode}
\usepackage{tikz}
\usetikzlibrary{arrows.meta, positioning, calc}
\usepackage{amsmath,amssymb}


\newcommand{\paperref}[1]{
\vfill
{\tiny\color{gray}#1}
}

% KU green colors from the original blog
\definecolor{kugreen}{RGB}{50,93,61}
\definecolor{kugreenlys}{RGB}{132,158,139}
\definecolor{kugreenlyslys}{RGB}{173,190,177}
\definecolor{kugreenlyslyslys}{RGB}{214,223,216}

\mode<presentation>
{
  \usetheme[numbers,totalnumber,compress,sidebarshades]{PaloAlto}
  \usecolortheme[named=kugreen]{structure}
  \useinnertheme{circles}
  \usefonttheme[onlymath]{serif}
  \setbeamercovered{transparent}
  \setbeamertemplate{blocks}[rounded][shadow=true]
  \setbeamertemplate{footline}[frame number]
}

% Optional logo, upload logo.pdf/png to Overleaf first
% \logo{\includegraphics[width=1.5cm]{logo}}

% empty short title [] -> nothing shown in the sidebar
\title[]{An Unreasonably Cheap Algorithm to Generate Extreme Weather Events}
\author[]{Alain Joss}
\institute{SDSC}
\date{\today}

% Section title slide before each new section
\AtBeginSection[]
{
  \begin{frame}[plain]
    \centering
    \vfill
    {\Huge \insertsection}
    \vfill
  \end{frame}
}
\begin{document}

\frame{\titlepage}

\section{Motivation}

\frame{
\frametitle{Critical systems and extreme weather}
\begin{itemize}
  \item Energy, transport, agriculture: weather dependent
  \begin{itemize}
    \item Short-to-medium-term operation $\leftarrow$ forecasting
    \item Long-term planning $\leftarrow$ scenario analysis (including extremes)
  \end{itemize}
  \item Extremes: unlikely but highly relevant for decision-making (potentially disruptive)
  \item Analyze their consequences $\rightarrow$ map bottlenecks \& points of failure
  \item Extreme scenarios must be evaluated before any forecast can see them
  \begin{itemize}
    \item some impacts are unavoidable $\rightarrow$ response decisions still have to be prepared ahead of the event
    \item black-swan type: arguably very highly relevant
  \end{itemize}
  \item Deciding whether, and how, to prepare for extreme weather involves trading off
        the cost of preparing for an event which might not occur, against the loss
        incurred if no preparations are made and it does occur
\end{itemize}
\vfill
$\rightarrow$ \textbf{How can these scenarios be generated?}
}

\frame{
\frametitle{Our objective}
\textbf{Objective:} simulate \emph{extreme}, yet \emph{plausible} weather scenarios
(in contrast to forecasting)
\begin{itemize}
  \item Scenarios: simulate weather at different intensity levels
  \item Example: if temperature goes up by $2^{\circ}$, \dots, $6^{\circ}$
        in this region of interest $\rightarrow$ then \dots
\end{itemize}
\vfill
\textbf{Definitions:}
\begin{itemize}
  \item \emph{Extreme}
  \begin{itemize}
    \item unlikely (and disruptive)
    \item lives on the tails of the climatological and forecasting distribution
  \end{itemize}
  \item \emph{Plausible}
  \begin{itemize}
    \item transitions are physically grounded
    \item variables co-vary in a coherent way
  \end{itemize}
\end{itemize}
\vfill
We give a precise description of the scenarios.
}

\frame{
\frametitle{Example: stress-testing power-grid control}
\begin{itemize}
  \item Climate changing fast $\rightarrow$ summers increasingly disruptive
  \item Goal: stress-test grid resilience ahead of the coming years
  \item Focus scenario: \textbf{heatwave}
  \begin{itemize}
    \item co-occurring extremes: heat, solar radiation, wind, fires, thunderstorms, hail
  \end{itemize}
  \item Potential consequences:
  \begin{itemize}
    \item grid instability
    \item infrastructure damage
    \item cascading failures
    \item blackouts
  \end{itemize}
\end{itemize}
}

\section{How to generate extreme weather events}

\frame{
\frametitle{Our approach}
\begin{itemize}
  \item \textbf{Model:}
  \begin{itemize}
    \item medium-range generative probabilistic weather model
    \item generative: score / diffusion / flow matching
  \end{itemize}
  \item \textbf{Guidance methods:}
  \begin{itemize}
    \item borrowed from the image generation domain
    \item translated to the weather domain
    \item two families: sampling-based or gradient-based
  \end{itemize}
\end{itemize}
}

\frame{
\frametitle{But \dots}
\textbf{Why not just use the tails of the base forecast distribution?}
\begin{itemize}
  \item Extremes $=$ rare tail events, possibly unprecedented
  \item May be impossible to obtain from the base forecasting model
  \item Or simply too expensive to sample
  \item We want a \emph{distribution} of extremes $\rightarrow$ even harder
  \item $\rightarrow$ need some scheme on top of the base model
\end{itemize}
\vfill
\textbf{Why use a generative model?}
\begin{itemize}
  \item The hope is to keep the generated extremes close to the learned data distribution
  \item If the modes of the learned data distribution encode physical realism, the
        generative model can help us achieve extremes that satisfy this property
\end{itemize}
}

\frame{
\frametitle{Images vs.\ extreme weather}
\textbf{Core challenge of our work:} translate / design guidance to match the
challenges of the weather domain.
\vfill
3 key differences:
\begin{itemize}
  \item \textbf{Steps:}
  \begin{itemize}
    \item image domain $\rightarrow$ final image obtained one-shot
    \item extreme-weather domain $\rightarrow$ multi-step trajectory
          (video-alike, but not really \dots)
  \end{itemize}
  \item \textbf{Evaluation:}
  \begin{itemize}
    \item image domain $\rightarrow$ visual appeal
    \item extreme-weather domain $\rightarrow$ visual checks $+$ quantified realism
  \end{itemize}
  \item \textbf{Target:}
  \begin{itemize}
    \item image domain $\rightarrow$ inherent to a specific task
          (super-resolution, inpainting, prompt alignment, \dots)
    \item extreme-weather domain $\rightarrow$ unclear a priori (designed)
  \end{itemize}
\end{itemize}
\vfill
Video-alike but not really $\rightarrow$ transitions must be physically grounded.
}

\frame{
\frametitle{Related work}
\begin{itemize}
  \item guide first step
  \item shift sampled noise
  \item plain classifier guidance
\end{itemize}
}

\section{Weather guidance}

\frame{
\frametitle{Data}
\textbf{ERA5 Reanalysis (ECMWF)}
\begin{itemize}
  \item Global, $1.5^{\circ}$ grid --- $240 \times 121$ points
  \item 82 fields per state:
  \begin{itemize}
    \item 4 surface: 2m temperature, MSL pressure, 10m u/v wind
    \item 6 upper-air $\times$ 13 pressure levels: geopotential, temperature,
          humidity, u/v wind, vertical velocity
  \end{itemize}
  \item 24-hour forecast steps, available at 00:00, 06:00, 12:00, 18:00 UTC
  \item Test set: $\geq$ 01.01.2020
\end{itemize}
}

\frame{
\frametitle{Model}
Prediction $=$ deterministic model $+$ residual flow matching model
\vfill
\centering
\resizebox{\textwidth}{!}{% Computational graph: deterministic core + generative residual flow.
% Self-contained tikzpicture (styles inline); requires the kugreen* colors and
% \usetikzlibrary{arrows.meta, positioning, calc}. Shared by presentation.tex
% (Model frame) and standalone.tex (PNG export).
\begin{tikzpicture}[
  font=\small,
  state/.style={draw=kugreen, fill=kugreenlyslyslys, rounded corners=2pt,
                minimum height=0.72cm, inner sep=4pt},
  model/.style={draw=kugreen, fill=kugreen, text=white, rounded corners=2pt,
                minimum height=0.9cm, inner sep=5pt},
  emb/.style={draw=kugreenlys, fill=kugreenlyslyslys!40, rounded corners=2pt,
              inner sep=3pt, font=\scriptsize},
  op/.style={draw=kugreen, circle, fill=white, inner sep=1pt},
  flow/.style={-{Latex[length=2.2mm]}, kugreen!80!black, thick},
  loop/.style={-{Latex[length=2.2mm]}, kugreen!80!black, thick, dashed},
  lbl/.style={font=\scriptsize, text=kugreen!60!black},
]

% ===== lane 1: deterministic pipeline (inputs on the same row) =====
\node[state] (xprev) at (0,3.0) {$x_{n-1}$};
\node[state] (xcur)  at (1.1,3.0) {$x_{n}$};
\node[op]    (catdet) at (2.05,3.0) {$\|$};
\node[model] (det)   at (3.9,3.0) {deterministic model};
\node[state] (xdet)  at (6.6,3.0) {$\hat{x}^{\text{det}}$};

\draw[flow] (xprev.north) to[out=55, in=125] (catdet.north);
\draw[flow] (xcur) -- (catdet);
\draw[flow] (catdet) -- (det);
\draw[flow] (det) -- (xdet);

% ===== time embeddings (month/hour shared by both models; e_t only for gen) =====
\node[emb] (temb) at (3.9,1.7) {$e_{\text{month}}+e_{\text{hour}}$};
\node[emb] (et)   at (5.6,1.7) {$e_t$};
\draw[flow] (temb) -- (det);

% ===== lane 2: residual flow (T Euler steps) =====
\node[state] (z0) at (-4.7,-0.8) {$z^{(0)}\!\sim\!\mathcal{N}(0,\sigma^{2}\mathbf{I})$};
\node[state] (zi) at (-2.3,-0.8) {$z^{(t)}$};
\node[state] (catgen) at (0.6,-0.8)
  {$\bigl[\, z^{(t)} \,\|\, x_{n-1} \,\|\, x_{n} \,\|\, \hat{x}^{\text{det}} \,\bigr]$};
\node[model] (gen) at (3.9,-0.8) {generative model};
\node[state] (ri) at (6.4,-0.8) {$r_t$};
\node[state] (ui) at (6.4,-2.5) {$u_t=\dfrac{r_t-z^{(t)}}{s_t}$};
\node[state] (euler) at (2.6,-2.5) {$z^{(t+1)}=z^{(t)}+h_t\,u_t$};

\draw[flow] (temb) -- (gen);
\draw[flow] (et.south) -- ($(gen.north)+(0.55,0)$);
\draw[flow] (z0) -- node[lbl, above] {$t=0$} (zi);
\draw[flow] (zi) -- (catgen);
\draw[flow] (catgen) -- (gen);
\draw[flow] (gen) -- (ri);
\draw[flow] (ri) -- (ui);
\draw[flow] (ui) -- (euler);
\draw[loop] (euler.west) -| node[lbl, pos=0.25, below] {repeat $t=0,\dots,T\!-\!1$} (zi.south);

% conditioning wires into the concatenation (the det output is re-used!)
\draw[flow] (xprev.south) to[out=-90, in=90] ($(catgen.north)+(-0.35,0)$);
\draw[flow] (xcur.south)  to[out=-90, in=90] ($(catgen.north)+(0.4,0)$);
\draw[flow, rounded corners=4pt] (xdet.south) |- (1.8,0.6)
  -- ($(catgen.north)+(1.2,0)$);

% ===== composition =====
\node[op]    (plus)   at (10.9,3.0) {$\oplus$};
\node[state] (xfinal) at (12.2,3.0) {$\hat{x}_{n+1}$};
\node[state] (zT)     at (10.9,-3.7) {$\sigma_r \odot z^{(T)}$};

\draw[flow] (euler.south) |- node[lbl, pos=0.75, below] {after $T$ steps} (zT.west);
\draw[flow] (zT.north) -- (plus.south);
\draw[flow] (xdet.east) -- (plus.west);
\draw[flow] (plus) -- (xfinal);

% caption line
\node[lbl, anchor=north] at (3.7,-4.5)
  {$s_t$: noise level \quad $h_t$: Euler step size \quad $\sigma_r$: residual scale
   \quad $\hat{x}_{n+1}=\hat{x}^{\text{det}}+\sigma_r \odot z^{(T)}$};

\end{tikzpicture}%
}
}

\frame{
\frametitle{Guidance @ weather time: pipeline}
We target the following:
\begin{itemize}
  \item heatwave scenarios (surface temperature)
  \item over a specific region
  \item with a specified change trajectory
  \item over a specific timeframe
  \item with respect to the unguided step
\end{itemize}
}

\frame{
\frametitle{Guidance @ weather time: mask}
\textbf{Mask:}
\begin{itemize}
  \item Gaussian over the region of interest
  \item at the variable of interest (surface temperature)
\end{itemize}
}

\frame{
\frametitle{Guidance @ weather time: change trajectory}
\textbf{Change trajectory:}
\begin{itemize}
  \item over an $N$-step timeframe
  \item defined as a percentage change trajectory
  \item with respect to the unguided step average over the mask
  \item (later: with respect to the unguided--guided step)
\end{itemize}
}

\frame{
\frametitle{Guidance @ flow time: pipeline}
At every flow step $t$, we steer the velocity before taking the Euler step:
\begin{enumerate}
  \item Estimate the clean state from the current noisy residual:
        $\hat{x}_t = \hat{x}^{\text{det}} + \sigma_r \odot \left( z^{(t)} + s_t\, u_t \right)$
  \item Score it with the masked loss: $L(\hat{x}_t)$
  \item Take the gradient w.r.t.\ the noisy state:
        $g_t = \nabla_{z^{(t)}} L(\hat{x}_t)$
  \item Steer the velocity: $\tilde{u}_t = u_t - \lambda_t\, g_t$
        \; with strength $\lambda_t = w \cdot a_t$
  \item Euler step as usual: $z^{(t+1)} = z^{(t)} + h_t\, \tilde{u}_t$
\end{enumerate}
\vfill
Two design choices define the method:
\begin{itemize}
  \item \textbf{Masked loss} --- what to push and where
  \item \textbf{Guidance strength--schedule} --- how hard and when
\end{itemize}
}

\frame{
\frametitle{Guidance @ flow time: probabilistic view (LBG)}

\textbf{Steering with a loss gradient $=$ sampling from a posterior:}

\[
p(x|y) \propto p(y|x)\,p(x)
\]
\[
\nabla_x\log p(x|y)
=
\nabla_x\log p(x)
+
\nabla_x\log p(y|x)
\]

\begin{block}{Guided vector field}
\[
\tilde{u}_t
=
u_t - \lambda_t\, g_t,
\qquad
g_t = \nabla_{z^{(t)}} L(\hat{x}_t) = -\nabla_{z^{(t)}}\log p(y|z^{(t)})
\]
\end{block}

\begin{itemize}
    \item $u_t$ transports the prior $p(x)$ --- the learned weather distribution
    \item the masked loss defines the likelihood: $p(y|x) \propto e^{-L}$
    \item exact posterior needs $w=1$ and a path-specific $a_t$
          (e.g.\ $\tfrac{s_t}{1-s_t}$ for linear paths)
    \item we trade exactness for \emph{control}: $\lambda_t = w \cdot a_t$
          with the gap-closing schedule
    \item unlike DPS, no rescaling of $w\,a_t$ by the gradient norm
\end{itemize}

\paperref{Dhariwal \& Nichol, NeurIPS 2021 \quad\textperiodcentered\quad Chung et al.\ (DPS), ICLR 2023}
}

\frame{
\frametitle{Guidance @ flow time: loss}
\begin{equation*}
  L = \Big( \textstyle\sum m \odot \hat{x}_t \;-\; (1+\delta) \sum m \odot x^{\text{ref}} \Big)^2
\end{equation*}
\begin{itemize}
  \item $\sum m \odot x$ --- the masked statistic: one number summarizing the region of interest
  \item $\hat{x}_t$ --- the clean-state estimate at flow step $t$ (where we currently are)
  \item $(1+\delta)\, \sum m \odot x^{\text{ref}}$ --- the target: the reference statistic, shifted by $\delta$
  \item $x^{\text{ref}}$ --- the model's own unguided forecast (or ground truth)
  \item the signed difference inside the square is the \textbf{gap} $r_t$ ---
        the quantity the schedule drives to zero
\end{itemize}
\vfill
The gradient flows \emph{through the model} into $z^{(t)}$:
the network Jacobian spreads the correction over the whole field.
}

\frame{
\frametitle{Guidance @ flow time: strength--schedule}
Applied strength factorizes: $\lambda_t = w \cdot a_t$
\begin{itemize}
  \item $w$ --- global gain: \emph{how hard} to push
  \item $a_t$ --- schedule: \emph{when} to push
\end{itemize}
\vfill
\textbf{Gap closing:} each step closes a fraction $\eta$ of the remaining gap:
\begin{equation*}
  r_{t+1} = (1-\eta)\, r_t
\end{equation*}
This prescribes the residual path $r_t = (1-\eta)^{t+1}\, r_0$
and the schedule $a_t = (1-\eta)^{t+1}$.
\vfill
Early steps do the heavy lifting; late steps only refine.
}

\frame{
\frametitle{Change trajectory}
\begin{itemize}
  \item We define the full change trajectory at all steps $n$
  \item The perturbation paper does it only at $n=1$, applies the forecasting model
        $N$ times, and lets the gradient flow back through all $N$ weather steps
  \item This enables us to control the full trajectory instead, and to have a
        distribution \emph{over the defined trajectory} --- compared to a distribution
        over possible ways to reach the final target state at $N$
\end{itemize}
}

\frame{
\frametitle{Guidance methods}
\begin{itemize}
  \item \textbf{OPTIMIZE-STRENGTH}
  \item \textbf{CLOSE-GAP}
\end{itemize}
}

\frame{
\frametitle{OPTIMIZE-STRENGTH}
Optimize the single scalar $w$ against the final loss --- no learning rate.
\begin{itemize}
  \item Run the guided flow with constant $w$: $\lambda_t = w \cdot a_t$
  \item $w$ is a scalar, so $\frac{dL}{dw}$ is an \emph{exact} derivative (one adjoint pass)
  \item Solve $\frac{dL}{dw} = 0$ with the secant method --- no learning rate,
        no iteration count; it stops when the final loss is closed,
        then reruns once with $w^{*}$
\end{itemize}
\vfill
Cost: a few full flow evaluations.
}

\frame{
\frametitle{CLOSE-GAP}
No $w$ at all --- close the gap \emph{exactly} at every step.
\begin{itemize}
  \item Measure the current gap $r_t$, aim at the prescribed path:
        $r_t^{\text{target}} = (1-\eta)^{t+1}\, r_0$
  \item The exact landing is a Newton move along the gradient:
        \begin{equation*}
          \lambda_t = \frac{2\, r_t \left( r_t - r_t^{\text{target}} \right)}{h_t\, \lVert g_t \rVert^2}
        \end{equation*}
  \item Anchored to the \emph{measured} gap $\rightarrow$ drift and overshoot are
        corrected at the next step
  \item $\eta = 1$: whole gap in one step; smaller $\eta$: spread smoothly over the flow
\end{itemize}
\vfill
Cost: a single guided flow --- no outer optimization.
}

\frame{
\frametitle{Notation summary}
% TODO: fill in
}

\frame{
\frametitle{Assumptions}
\begin{itemize}
  \item If we can find metrics that enable us to make statements about the
        plausibility of the weather states, we don't need to make any assumptions
        about the physical realism of the guidance target and the learned model function
  \item If this holds, we can also make statements about full trajectories:
        if a one-step guided forecast is physically plausible, then so is its
        autoregressive successor --- or better, we can measure the degradation
        of physical realism
  \item Reason to guide this way: many little steps that only slightly change
        the trajectory are good
\end{itemize}
}

\section{Experimental setup}

\frame{
\frametitle{Validation}
Guiding towards ground truth
}

\section{Results}

\frame{
\frametitle{Results}
% TODO: fill in
}

\section{Limitations}

\frame{
\frametitle{Limitations}
\begin{itemize}
  \item Works well for heatwaves because the mask is fixed $\rightarrow$
        not appropriate for cyclones and other moving events
  \item Sampling methods provide a more direct way to respect the transition function
        learned by the model $\rightarrow$ no manual mask and manual trajectories
        $\rightarrow$ would side-step the need to check the realism of both mask and
        delta trajectory, and of the produced results
\end{itemize}
}

\section{Future directions}

\frame{
\frametitle{Future directions}
\begin{itemize}
  \item GenCast: higher resolution --- $0.25^{\circ}$ $==$ $721 \times 1440$ grid
  \item Likelihood estimation of an event
  \item Guiding noise
  \item Guiding steps
\end{itemize}
}

\appendix
\section{Appendix}

\frame{
\frametitle{Residual model}
\textbf{Inputs:}
\begin{itemize}
  \item \texttt{state} --- the current analysis $x_0$ (at $t_0$)
  \item \texttt{prev\_state} --- the state at $t_0 - 24\,$h
  \item \texttt{pred\_state} $=$ \texttt{det\_pred} --- the deterministic prediction
  \item $z_t$ --- the noisy residual latent at the current noise level
        (the thing being denoised)
  \item time embedding --- month $+$ hour $+$ flow timestep $t$:
        seasonality/diurnal conditioning $+$ noise-level conditioning
\end{itemize}
}

\frame{
\frametitle{Anchoring scenarios in time}
\textbf{Option 1: take past events}
\begin{itemize}
  \item extremify them
  \item compare against available ground truth
\end{itemize}
\vfill
\textbf{Option 2: generate from the present into the future}
\begin{itemize}
  \item the past may not contain all scenarios of interest
  \item stress-test infrastructure in its current state
  \item ground truth not available
\end{itemize}
}

\end{document}
